\documentclass[11pt]{article}

\usepackage{amsmath,amssymb}
\usepackage{xurl}
\usepackage[colorlinks=true,linkcolor=blue,citecolor=blue,urlcolor=blue]{hyperref}

% Shared commands for notes in docs/paper-gaps/.

\DeclareUnicodeCharacter{2080}{\ifmmode{}_{0}\else\textsubscript{0}\fi}
\DeclareUnicodeCharacter{2081}{\ifmmode{}_{1}\else\textsubscript{1}\fi}
\DeclareUnicodeCharacter{2082}{\ifmmode{}_{2}\else\textsubscript{2}\fi}
\DeclareUnicodeCharacter{2099}{\ifmmode{}_{n}\else\textsubscript{n}\fi}

\newcommand{\Lean}{\textsc{Lean}}
\ExplSyntaxOn
\NewDocumentCommand{\leanid}{m}{
  \ifmmode
    \textsf{\detokenize{#1}}
  \else
    \group_begin:
    \urlstyle{sf}
    \tl_set:Nn \l_tmpa_tl {#1}
    \tl_replace_all:Nnn \l_tmpa_tl {\_} {_}
    \exp_args:NV \path \l_tmpa_tl
    \group_end:
  \fi
}
\ExplSyntaxOff
\providecommand{\tr}{\operatorname{tr}}
\newcommand{\tnleanIssueBase}{https://github.com/LionSR/TNLean/issues/}
\newcommand{\tnleanPrBase}{https://github.com/LionSR/TNLean/pull/}
\newcommand{\tnleanDocBase}{https://sirui-lu.com/TNLean/docs/}
\newcommand{\ghissue}[1]{\href{\tnleanIssueBase#1}{GitHub issue \##1}}
\newcommand{\ghpr}[1]{\href{\tnleanPrBase#1}{PR \##1}}
\newcommand{\doclink}[1]{\href{\tnleanDocBase#1.html}{\path{#1}}}

% Dirac notation and the identity operator, as in blueprint/src/macros/common.tex.
% \providecommand keeps note-local definitions authoritative.
\usepackage{dsfont}
\providecommand{\ket}[1]{|#1\rangle}
\providecommand{\bra}[1]{\langle #1|}
\providecommand{\braket}[2]{\langle #1 | #2 \rangle}
\providecommand{\ketbra}[2]{|#1\rangle\!\langle #2|}
\providecommand{\Id}{\mathds{1}}
\providecommand{\one}{\mathds{1}}
\providecommand{\C}{\mathbb{C}}
\providecommand{\MN}[1]{M_{#1}(\C)}
\providecommand{\MD}{M_D(\C)}

% External referencing for paper sources.  The build helper writes sanitized
% auxiliary files under build/paper-aux/ and excludes duplicated source labels.
\usepackage{xr}

% Import a generated paper auxiliary file with a caller-supplied label prefix.
% The candidate paths support builds from docs/paper-gaps/, the repository root,
% and build/paper-aux/ itself.
\newcommand{\importpaperaux}[2]{%
  \IfFileExists{../../build/paper-aux/#2.aux}{%
    \externaldocument[#1]{../../build/paper-aux/#2}%
  }{%
    \IfFileExists{build/paper-aux/#2.aux}{%
      \externaldocument[#1]{build/paper-aux/#2}%
    }{%
      \IfFileExists{#2.aux}{%
        \externaldocument[#1]{#2}%
      }{}%
    }%
  }%
}

% General external reference macros
\newcommand{\paperref}[2]{\ref{#1-#2}}
\newcommand{\papereqref}[2]{(\ref{#1-#2})}

% CPSV16 source (1606.00608 / MPDO-22-12-17-2.tex) external document and shortcuts
\importpaperaux{cpsv16-}{1606.00608/MPDO-22-12-17-2-tnlean}
\newcommand{\cpsvref}[1]{\paperref{cpsv16}{#1}}
\newcommand{\cpsveqref}[1]{\papereqref{cpsv16}{#1}}

% Verdict marker: \gapnote{<kind>}{<status>}. Kind is the policy
% classification (clarification, local-correction, scope-restriction,
% unfaithful, false-source, open-gap); status is open, wip (elimination
% underway), resolved, or historical. Machine-read by the site index;
% prints a verdict line.
\newcommand{\gapnote}[2]{\par\noindent\textbf{Verdict:} #1 (#2).\par}


\title{Complete-Positivity Scope of the Peripheral Cyclicity Theorem}
\author{}
\date{\today}

\begin{document}
\maketitle

\section{The source theorem}

Theorem 6.6 of \emph{Quantum Channels \& Operations} (Wolf 2012), stated at
lines 244--252 of \texttt{Notes/WolfNotePDF/wolf\_ch6\_sections.txt}, concerns
an irreducible positive unital map $T:M_D(\mathbb C)\to M_D(\mathbb C)$ that
satisfies the Schwarz inequality. Item~(1) concludes that the peripheral
spectrum is a finite cyclic group of order $m\leq D^2$; items~(2)--(4) give
nondegeneracy, a peripheral unitary, and cyclic spectral projections.
Complete positivity and a Kraus representation are not hypotheses of the
theorem.

\section{The restriction in the formalization}

The available Lean statement is
\leanid{MPSTensor.peripheralEigenvalues_eq_range_primitiveRoot} in
\texttt{TNLean/Channel/Peripheral/CyclicGroup.lean}. It is stated for a finite
matrix family $K$ and its Kraus map
\[
    E(X)=\sum_i K_iXK_i^\dagger.
\]
It assumes that this Kraus map is unital and irreducible and that its adjoint
has a positive-definite fixed point. A finite Kraus presentation implies
complete positivity, which is strictly stronger than Wolf's positive Schwarz
hypothesis. Thus the Lean theorem proves the completely positive
specialization of the cyclicity conclusion in Wolf Theorem 6.6(1), but not its
bound $m\leq D^2$ and not items~(2)--(4). The corresponding Blueprint entry
states this restriction explicitly and uses its own specialization title.

\section{Elimination plan}

To remove the restriction, formulate peripheral eigenvectors and their
multiplicative-domain relations for an abstract positive unital Schwarz map.
The proof must establish invertibility of peripheral eigenvectors and closure
of peripheral eigenvalues under multiplication without choosing Kraus
operators. The finite-group argument in the present cyclicity theorem then
applies unchanged. Once that abstract result exists, the Kraus theorem should
be retained as a specialization rather than serving as the source-facing form
of Wolf Theorem 6.6.

\section{Verdict}

\textbf{Verdict: scope restriction.} The formal theorem proves only the
cyclicity conclusion of Wolf Theorem 6.6(1), under the stronger
complete-positivity hypothesis supplied by a finite Kraus representation. It
does not include the bound $m\leq D^2$ or items~(2)--(4). The abstract
Schwarz-map cyclicity statement remains open.

\end{document}
